\title{2021 年高考数学 全国甲卷-理科}
\date{}
\maketitle

\section*{一、单项选择题}
本题共12小题，每小题5分，共60分．在每小题给出的四个选项中，只有一项是符合题目要求的．

\begin{question}[points=5]
设集合 $M = \{x \mid 0 < x < 4\}$, $N = \left\{x \mid \frac{1}{3} \le x \le 5\right\}$, 则 $M \cap N =$ ( )
\begin{choices}
  \item $\left\{x \mid 0 < x \le \frac{1}{3}\right\}$
  \item $\left\{x \mid \frac{1}{3} \le x < 4\right\}$
  \item $\{x \mid 4 \le x < 5\}$
  \item $\{x \mid 0 < x \le 5\}$
\end{choices}
\begin{solution}
\textbf{答案：}B

\textbf{分析：}根据交集定义运算即可

\textbf{详解：}$M = \{x \mid 0 < x < 4\}$, $N = \left\{x \mid \frac{1}{3} \le x \le 5\right\}$, 则 $M \cap N = \left\{x \mid \frac{1}{3} \le x < 4\right\}$．
\end{solution}
\end{question}

\begin{question}[points=5]
为了解某地农村经济情况，对该地农户家庭年收入进行抽样调查，将农户家庭年收入的调查数据整理得到如下频率分布直方图：根据此频率分布直方图，下面结论中不正确的是（ ）

\begin{center}
  \includegraphics[width=.38\linewidth]{question-2-1.jpg}
\end{center}
\begin{choices}
  \item 该地农户家庭年收入低于4.5万元的农户比率估计为6%
  \item 该地农户家庭年收入不低于10.5万元的农户比率估计为10%
  \item 估计该地农户家庭年收入的平均值不超过6.5万元
  \item 估计该地有一半以上的农户，其家庭年收入介于4.5万元至8.5万元之间
\end{choices}
\begin{solution}
\textbf{答案：}C

\textbf{分析：}根据直方图的意义直接计算相应范围内的频率，即可判定ABD，以各组的中间值作为代表乘以相应的频率，然后求和即得到样本的平均数的估计值，也就是总体平均值的估计值，计算后即可判定C．

\textbf{详解：}由频率分布直方图可得：$0.02+0.04=0.06=6\%$，故A正确；$0.04+0.02\times3=0.10=10\%$，故B正确；$0.10+0.14+0.20\times2=0.64=64\%>50\%$，故D正确；平均值估计为$3\times0.02+4\times0.04+5\times0.10+6\times0.14+7\times0.20+8\times0.20+9\times0.10+10\times0.10+11\times0.04+12\times0.02+13\times0.02+14\times0.02=7.68$万元，超过6.5万元，故C错误．
\end{solution}
\end{question}

\begin{question}[points=5]
已知 $(1-i)^2 z = 3 + 2i$, 则 $z =$ ( )
\begin{choices}
  \item $-1-\frac{3}{2}i$
  \item $-1+\frac{3}{2}i$
  \item $-\frac{3}{2}+i$
  \item $-\frac{3}{2}-i$
\end{choices}
\begin{solution}
\textbf{答案：}B

\textbf{分析：}由已知得 $z = \frac{3+2i}{-2i}$，根据复数除法运算法则，即可求解．

\textbf{详解：}$(1-i)^2 = -2i$, 则 $z = \frac{3+2i}{-2i} = \frac{(3+2i)\cdot i}{-2i\cdot i} = \frac{-2+3i}{2} = -1+\frac{3}{2}i$．
\end{solution}
\end{question}

\begin{question}[points=5]
青少年视力是社会普遍关注的问题，视力情况可借助视力表测量．通常用五分记录法和小数记录法记录视力数据，五分记录法的数据 $L$ 和小数记录表的数据 $V$ 满足 $L = 5 + \lg V$．已知某同学视力的五分记录法的数据为 $4.9$，则其视力的小数记录法的数据为（ ）($\sqrt[10]{10} \approx 1.259$)
\begin{choices}
  \item $1.5$
  \item $1.2$
  \item $0.8$
  \item $0.6$
\end{choices}
\begin{solution}
\textbf{答案：}C

\textbf{分析：}根据 $L,V$ 关系，当 $L=4.9$ 时，求出 $\lg V$，再用指数表示 $V$，即可求解．

\textbf{详解：}由 $L=5+\lg V$，当 $L=4.9$ 时，$\lg V = -0.1$，则 $V = 10^{-0.1} = 10^{-\frac{1}{10}} = \frac{1}{\sqrt[10]{10}} \approx \frac{1}{1.259} \approx 0.8$．
\end{solution}
\end{question}

\begin{question}[points=5]
已知 $F_1, F_2$ 是双曲线 $C$ 的两个焦点，$P$ 为 $C$ 上一点，且 $\angle F_1PF_2 = 60^\circ$, $|PF_1| = 3|PF_2|$，则 $C$ 的离心率为（ ）
\begin{choices}
  \item $\frac{\sqrt{7}}{2}$
  \item $\frac{\sqrt{13}}{2}$
  \item $\sqrt{7}$
  \item $\sqrt{13}$
\end{choices}
\begin{solution}
\textbf{答案：}A

\textbf{分析：}根据双曲线的定义及条件，表示出 $PF_1, PF_2$，结合余弦定理可得答案．

\textbf{详解：}由双曲线的定义可得 $|PF_1|-|PF_2| = 2|PF_2| = 2a$, 所以 $|PF_2| = a$, $|PF_1| = 3a$；在 $\triangle F_1PF_2$ 中，由余弦定理得 $4c^2 = 9a^2 + a^2 - 2\cdot 3a\cdot a\cos60^\circ$, 解得 $4c^2 = 7a^2$, 所以 $e^2 = \frac{c^2}{a^2} = \frac{7}{4}$, 即 $e = \frac{\sqrt{7}}{2}$．
\end{solution}
\end{question}

\begin{question}[points=5]
在一个正方体中，过顶点 $A$ 的三条棱的中点分别为 $E, F, G$．该正方体截去三棱锥 $A-EFG$ 后，所得多面体的三视图中，正视图如图所示，则相应的侧视图是（ ）

\begin{center}
  \includegraphics[width=.38\linewidth]{question-6-1.jpg}
\end{center}
\begin{choices}
  \item 
  \item 
  \item 
  \item 
\end{choices}
\begin{solution}
\textbf{答案：}D

\textbf{分析：}根据题意及题目所给的正视图还原出几何体的直观图，结合直观图进行判断．

\textbf{详解：}由题意及正视图可得几何体的直观图，如图所示，所以其侧视图为D．
\end{solution}
\end{question}

\begin{question}[points=5]
等比数列 $\{a_n\}$ 的公比为 $q$，前 $n$ 项和为 $S_n$．设甲：$q>0$，乙：$\{S_n\}$ 是递增数列，则（ ）
\begin{choices}
  \item 甲是乙的充分条件但不是必要条件
  \item 甲是乙的必要条件但不是充分条件
  \item 甲是乙的充要条件
  \item 甲既不是乙的充分条件也不是乙的必要条件
\end{choices}
\begin{solution}
\textbf{答案：}B

\textbf{分析：}当 $q>0$ 时，通过举反例说明甲不是乙的充分条件；当 $\{S_n\}$ 是递增数列时，必有 $a_n>0$ 成立即可说明 $q>0$ 成立，则甲是乙的必要条件．

\textbf{详解：}当 $a_n = -2, -4, -8, \cdots$ 时，$q>0$，但 $\{S_n\}$ 递减，故甲不是乙的充分条件；若 $\{S_n\}$ 递增，则必有 $a_n>0$，进而 $q>0$，故甲是乙的必要条件．故选B．
\end{solution}
\end{question}

\begin{question}[points=5]
2020年12月8日，中国和尼泊尔联合公布珠穆朗玛峰最新高程为8848.86（单位：m），三角高程测量法是珠峰高程测量方法之一．如图是三角高程测量法的一个示意图，现有 $A, B, C$ 三点，且 $A, B, C$ 在同一水平面上的投影 $A', B', C'$ 满足 $\angle A'C'B' = 45^\circ$, $\angle A'B'C' = 60^\circ$．由 $C$ 点测得 $B$ 点的仰角为 $15^\circ$, $BB'$ 与 $CC'$ 的差为100；由 $B$ 点测得 $A$ 点的仰角为 $45^\circ$，则 $A, C$ 两点到水平面 $A'B'C'$ 的高度差 $AA' - CC'$ 约为（$\sqrt{3} \approx 1.732$）（ ）

\begin{center}
  \includegraphics[width=.38\linewidth]{question-8-1.jpg}
\end{center}
\begin{choices}
  \item $346$
  \item $373$
  \item $446$
  \item $473$
\end{choices}
\begin{solution}
\textbf{答案：}B

\textbf{分析：}通过做辅助线，将已知所求量转化到一个三角形中，借助正弦定理，求得 $A'B'$，进而得到答案．

\textbf{详解：}过 $C$ 作 $CH \perp BB'$，过 $B$ 作 $BD \perp AA'$，则 $AA' - CC' = AA' - (BB' - BH) = AA' - BB' + 100 = AD + 100$．由题易知 $\triangle ADB$ 为等腰直角三角形，所以 $AD = DB$，故 $AA' - CC' = DB + 100 = A'B' + 100$．因为 $\angle BCH = 15^\circ$，所以 $CH = C'B' = \frac{100}{\tan 15^\circ}$．在 $\triangle A'B'C'$ 中，由正弦定理得 $\frac{A'B'}{\sin 45^\circ} = \frac{C'B'}{\sin 75^\circ} = \frac{100}{\tan 15^\circ \cos 15^\circ \sin 15^\circ}$，而 $\sin 15^\circ = \frac{\sqrt{6}-\sqrt{2}}{4}$，得 $A'B' = 100(\sqrt{3}+1) \approx 273$，所以 $AA' - CC' \approx 373$．
\end{solution}
\end{question}

\begin{question}[points=5]
若 $\alpha \in \left(0, \frac{\pi}{2}\right)$, $\tan 2\alpha = \frac{\cos \alpha}{2 - \sin \alpha}$, 则 $\tan \alpha = $ ( )
\begin{choices}
  \item $\frac{\sqrt{15}}{15}$
  \item $\frac{\sqrt{5}}{5}$
  \item $\frac{\sqrt{5}}{3}$
  \item $\frac{\sqrt{15}}{3}$
\end{choices}
\begin{solution}
\textbf{答案：}A

\textbf{分析：}由二倍角公式可得 $\tan 2\alpha = \frac{\sin 2\alpha}{\cos 2\alpha} = \frac{2\sin \alpha \cos \alpha}{1-2\sin^2 \alpha}$，再结合已知可求得 $\sin \alpha = \frac{1}{4}$，利用同角三角函数的基本关系即可求解．

\textbf{详解：}由条件得 $\frac{2\sin \alpha \cos \alpha}{1-2\sin^2 \alpha} = \frac{\cos \alpha}{2-\sin \alpha}$，因为 $\cos \alpha \neq 0$，所以 $\frac{2\sin \alpha}{1-2\sin^2 \alpha} = \frac{1}{2-\sin \alpha}$，解得 $\sin \alpha = \frac{1}{4}$，则 $\cos \alpha = \frac{\sqrt{15}}{4}$，故 $\tan \alpha = \frac{\sqrt{15}}{15}$．
\end{solution}
\end{question}

\begin{question}[points=5]
将4个1和2个0随机排成一行，则2个0不相邻的概率为（ ）
\begin{choices}
  \item $\frac{1}{3}$
  \item $\frac{2}{5}$
  \item $\frac{2}{3}$
  \item $\frac{4}{5}$
\end{choices}
\begin{solution}
\textbf{答案：}C

\textbf{分析：}采用插空法，4个1产生5个空，分2个0相邻和2个0不相邻进行求解．

\textbf{详解：}4个1和2个0随机排成一行，总情况数为 $C_6^2 = 15$（或 $\frac{6!}{4!2!} = 15$）．2个0不相邻的排法：先排4个1有1种，再将2个0插入5个空，有 $C_5^2 = 10$ 种，所以概率为 $\frac{10}{15} = \frac{2}{3}$．
\end{solution}
\end{question}

\begin{question}[points=5]
已知 $A, B, C$ 是半径为1的球 $O$ 的球面上的三个点，且 $AC \perp BC$, $AC = BC = 1$，则三棱锥 $O-ABC$ 的体积为（ ）
\begin{choices}
  \item $\frac{\sqrt{2}}{12}$
  \item $\frac{\sqrt{3}}{12}$
  \item $\frac{\sqrt{2}}{4}$
  \item $\frac{\sqrt{3}}{4}$
\end{choices}
\begin{solution}
\textbf{答案：}A

\textbf{分析：}由题可得 $\triangle ABC$ 为等腰直角三角形，得出 $\triangle ABC$ 外接圆的半径，则可求得 $O$ 到平面 $ABC$ 的距离，进而求得体积．

\textbf{详解：}因为 $AC \perp BC$, $AC = BC = 1$，所以 $AB = \sqrt{2}$，$\triangle ABC$ 外接圆半径 $r = \frac{\sqrt{2}}{2}$．球半径 $R=1$，则 $O$ 到平面 $ABC$ 的距离 $d = \sqrt{1^2 - \left(\frac{\sqrt{2}}{2}\right)^2} = \frac{\sqrt{2}}{2}$，故 $V_{O-ABC} = \frac{1}{3} \cdot \frac{1}{2} \cdot 1 \cdot 1 \cdot \frac{\sqrt{2}}{2} = \frac{\sqrt{2}}{12}$．
\end{solution}
\end{question}

\begin{question}[points=5]
设函数 $f(x)$ 的定义域为 $\mathbb{R}$，$f(x+1)$ 为奇函数，$f(x+2)$ 为偶函数，当 $x \in [1,2]$ 时，$f(x) = ax^2 + b$．若 $f(0) + f(3) = 6$，则 $f\left(\frac{9}{2}\right) =$ ( )
\begin{choices}
  \item $-\frac{9}{4}$
  \item $-\frac{3}{2}$
  \item $\frac{7}{4}$
  \item $\frac{5}{2}$
\end{choices}
\begin{solution}
\textbf{答案：}D

\textbf{分析：}通过 $f(x+1)$ 是奇函数和 $f(x+2)$ 是偶函数条件，可以确定出函数解析式 $f(x) = -2x^2 + 2$，进而利用定义或周期性结论，即可得到答案．

\textbf{详解：}由 $f(x+1)$ 为奇函数得 $f(-x+1) = -f(x+1)$，由 $f(x+2)$ 为偶函数得 $f(x+2) = f(-x+2)$．令 $x=1$ 得 $f(0) = -f(2) = -(4a+b)$，$f(3) = f(1) = a+b$，代入 $f(0)+f(3)=6$ 得 $-(4a+b)+a+b=6$，解得 $a=-2$．令 $x=0$ 得 $f(1)=-f(1)$，故 $f(1)=0$，所以 $b=2$，则 $f(x) = -2x^2+2$．由对称性可得周期 $T=4$，故 $f\left(\frac{9}{2}\right) = f\left(\frac{1}{2}\right) = -f\left(\frac{3}{2}\right)$，而 $f\left(\frac{3}{2}\right) = -2\cdot\frac{9}{4}+2 = -\frac{5}{2}$，所以 $f\left(\frac{9}{2}\right) = \frac{5}{2}$．
\end{solution}
\end{question}

\section*{二、填空题}
本题共4小题，每小题5分，共20分．

\begin{question}[points=5]
曲线 $y = \frac{2x-1}{x+2}$ 在点 $(-1, -3)$ 处的切线方程为．
\fillin[{$5x - y + 2 = 0$}]
\begin{solution}
\textbf{答案：}$5x - y + 2 = 0$

\textbf{分析：}先验证点在曲线上，再求导，代入切线方程公式即可．

\textbf{详解：}$y' = \frac{2(x+2)- (2x-1)}{(x+2)^2} = \frac{5}{(x+2)^2}$，则 $y'|_{x=-1} = 5$，所以切线方程为 $y+3=5(x+1)$，即 $5x - y + 2 = 0$．
\end{solution}
\end{question}

\begin{question}[points=5]
已知向量 $\vec{a} = (3,1)$, $\vec{b} = (1,0)$, $\vec{c} = \vec{a} + k\vec{b}$．若 $\vec{a} \perp \vec{c}$，则 $k = $ ．
\fillin[{$-\frac{10}{3}$}]
\begin{solution}
\textbf{答案：}$-\frac{10}{3}$

\textbf{分析：}利用向量的坐标运算法则求得向量 $\vec{c}$ 的坐标，利用向量的数量积为零求得 $k$ 的值．

\textbf{详解：}$\vec{c} = (3+k, 1)$，由 $\vec{a} \perp \vec{c}$ 得 $3(3+k) + 1\times 1 = 0$，解得 $k = -\frac{10}{3}$．
\end{solution}
\end{question}

\begin{question}[points=5]
已知 $F_1, F_2$ 为椭圆 $C: \frac{x^2}{16} + \frac{y^2}{4} = 1$ 的两个焦点，$P, Q$ 为 $C$ 上关于坐标原点对称的两点，且 $|PQ| = |F_1F_2|$，则四边形 $PF_1QF_2$ 的面积为．
\fillin[{$8$}]
\begin{solution}
\textbf{答案：}$8$

\textbf{分析：}根据已知可得 $PF_1 \perp PF_2$，设 $|PF_1| = m, |PF_2| = n$，利用勾股定理结合 $m+n=8$，求出 $mn$，四边形 $PF_1QF_2$ 面积等于 $mn$．

\textbf{详解：}由条件知 $|PQ| = |F_1F_2|$，所以四边形 $PF_1QF_2$ 为矩形．设 $|PF_1| = m, |PF_2| = n$，则 $m+n = 2a = 8$，$m^2+n^2 = 4c^2 = 4(16-4) = 48$，由 $(m+n)^2 = m^2+n^2+2mn$ 得 $64 = 48+2mn$，所以 $mn = 8$，即四边形 $PF_1QF_2$ 面积为 $8$．
\end{solution}
\end{question}

\begin{question}[points=5]
已知函数 $f(x) = 2\cos(\omega x + \varphi)$ 的部分图像如图所示，则满足条件 $\left(f(x) - f\left(-\frac{7\pi}{4}\right)\right) \left(f(x) - f\left(\frac{4\pi}{3}\right)\right) > 0$ 的最小正整数 $x$ 为．

\begin{center}
  \includegraphics[width=.38\linewidth]{question-16-1.jpg}
\end{center}
\fillin[{$2$}]
\begin{solution}
\textbf{答案：}$2$

\textbf{分析：}先根据图象求出函数 $f(x)$ 的解析式，再求出 $f(-\frac{7\pi}{4}), f(\frac{4\pi}{3})$ 的值，然后求解三角不等式可得最小正整数或验证数值可得．

\textbf{详解：}由图可知 $\frac{3}{4}T = \frac{13\pi}{12} - \frac{\pi}{3} = \frac{3\pi}{4}$，所以 $T = \pi$，$\omega = 2$；由五点法得 $2\cdot\frac{\pi}{3} + \varphi = \frac{\pi}{2}$，所以 $\varphi = -\frac{\pi}{6}$，故 $f(x) = 2\cos\left(2x - \frac{\pi}{6}\right)$．计算得 $f\left(-\frac{7\pi}{4}\right) = 2\cos\left(-\frac{11\pi}{3}\right) = 2\cos\frac{\pi}{3} = 1$，$f\left(\frac{4\pi}{3}\right) = 2\cos\left(\frac{5\pi}{2}\right) = 0$，所以不等式化为 $(f(x)-1)f(x) > 0$，即 $f(x) > 1$ 或 $f(x) < 0$．通过验证最小正整数 $x=1$ 时，$f(1) = 2\cos\left(2 - \frac{\pi}{6}\right) < 1$；$x=2$ 时，$f(2) = 2\cos\left(4 - \frac{\pi}{6}\right) < 0$，故最小正整数 $x=2$．
\end{solution}
\end{question}

\section*{三、解答题}
共70分．解答应写出文字说明、证明过程或演算步骤，第17～21题为必考题，每个试题考生都必须作答．第22、23题为选考题，考生根据要求作答．

\begin{question}[points=12]
甲、乙两台机床生产同种产品，产品按质量分为一级品和二级品，为了比较两台机床产品的质量，分别用两台机床各生产了200件产品，产品的质量情况统计如下表：\begin{tabular}{|c|c|c|c|} \hline & 一级品 & 二级品 & 合计 \\ \hline 甲机床 & 150 & 50 & 200 \\ \hline 乙机床 & 120 & 80 & 200 \\ \hline 合计 & 270 & 130 & 400 \\ \hline \end{tabular}（1）甲机床、乙机床生产的产品中一级品的频率分别是多少？（2）能否有99%的把握认为甲机床的产品质量与乙机床的产品质量有差异？附：$K^2 = \frac{n(ad-bc)^2}{(a+b)(c+d)(a+c)(b+d)}$，$P(K^2 \ge k)$：0.050、0.010、0.001对应 $k$：3.841、6.635、10.828．
\begin{solution}
\textbf{答案：}（1）75%；60%；（2）能．

\textbf{分析：}本题考查频率统计和独立性检验，根据给出公式计算即可．

\textbf{详解：}（1）甲机床一级品频率为 $\frac{150}{200}=75\%$，乙机床一级品频率为 $\frac{120}{200}=60\%$．（2）$K^2 = \frac{400\times(150\times80-50\times120)^2}{270\times130\times200\times200} = \frac{400}{39} \approx 10.256 > 6.635$，所以有99%的把握认为甲机床的产品质量与乙机床有差异．
\end{solution}
\end{question}

\begin{question}[points=12]
已知数列 $\{a_n\}$ 的各项均为正数，记 $S_n$ 为 $\{a_n\}$ 的前 $n$ 项和，从下面①②③中选取两个作为条件，证明另外一个成立．①数列 $\{a_n\}$ 是等差数列；②数列 $\{\sqrt{S_n}\}$ 是等差数列；③ $a_2 = 3a_1$．注：若选择不同的组合分别解答，则按第一个解答计分．
\begin{solution}
\textbf{答案：}答案见解析

\textbf{分析：}选①②作条件证明③时，可设 $\sqrt{S_n} = an+b$，结合 $a_n, S_n$ 的关系求出 $a_n$，利用 $\{a_n\}$ 是等差数列可证 $a_2=3a_1$；选①③作条件证明②时，根据等差数列的求和公式表示出 $S_n$，结合等差数列定义可证；选②③作条件证明①时，设出 $\sqrt{S_n}=an+b$，结合 $a_n, S_n$ 的关系求出 $a_n$，根据 $a_2=3a_1$ 可求 $b$，然后可证 $\{a_n\}$ 是等差数列．

\textbf{详解：}（1）选①②作条件证明③：设 $\sqrt{S_n}=an+b \ (a>0)$，则 $S_n = (an+b)^2$，$a_1 = (a+b)^2$，$n\ge2$ 时 $a_n = S_n - S_{n-1} = a(2an-a+2b)$．因为 $\{a_n\}$ 等差，所以 $\{a_n\}$ 通项为一次式，故 $a+b=0$？由 $a_2=3a_1$ 可推得 $b=0$，则 $a_n = a^2(2n-1)$，故 $a_2=3a_1$．（2）选①③作条件证明②：由 $a_2=3a_1$ 等差得公差 $d=2a_1$，$S_n = na_1 + \frac{n(n-1)}{2}\cdot 2a_1 = n^2 a_1$，所以 $\sqrt{S_n} = \sqrt{a_1} n$，是等差数列．（3）选②③作条件证明①：设 $\sqrt{S_n}=an+b$，同（1）得 $a_n = a(2an - a + 2b)$，由 $a_2=3a_1$ 得 $b=0$ 或 $b=-\frac{4a}{3}$，其中 $b=0$ 时 $a_n = a^2(2n-1)$ 为等差；$b=-\frac{4a}{3}$ 时 $S_1<0$ 不合题意，舍去．所以 $\{a_n\}$ 为等差数列．
\end{solution}
\end{question}

\begin{question}[points=12]
已知直三棱柱 $ABC-A_1B_1C_1$ 中，侧面 $AA_1B_1B$ 为正方形，$AB=BC=2$，$E, F$ 分别为 $AC$ 和 $CC_1$ 的中点，$D$ 为棱 $A_1B_1$ 上的点，$BF \perp A_1B_1$．（1）证明：$BF \perp DE$；（2）当 $B_1D$ 为何值时，面 $BB_1C_1C$ 与面 $DFE$ 所成的二面角的正弦值最小？
\begin{solution}
\textbf{答案：}（1）见解析；（2）$B_1D = \frac{1}{2}$

\textbf{分析：}通过已知条件，确定三条互相垂直的直线，建立合适的空间直角坐标系，借助空间向量证明线线垂直和求出二面角的平面角的余弦值最大，进而可以确定出答案．

\textbf{详解：}（1）因为直三棱柱，$BB_1 \perp$ 底面 $ABC$，所以 $BB_1 \perp AB$．又 $A_1B_1 \parallel AB$，$BF \perp A_1B_1$，所以 $BF \perp AB$，故 $AB \perp$ 平面 $BCC_1B_1$．以 $B$ 为原点，$BA, BC, BB_1$ 为 $x, y, z$ 轴建立空间直角坐标系，则 $B(0,0,0), A(2,0,0), C(0,2,0), B_1(0,0,2), A_1(2,0,2), C_1(0,2,2), E(1,1,0), F(0,2,1)$．设 $D(a,0,2)$（$0\le a\le 2$），则 $\overrightarrow{BF}=(0,2,1)$，$\overrightarrow{DE}=(1-a,1,-2)$，$\overrightarrow{BF}\cdot\overrightarrow{DE}=0$，故 $BF \perp DE$．（2）平面 $DFE$ 的法向量 $\vec{m}=(3,1+a,2-a)$，平面 $BCC_1B_1$ 的法向量 $\overrightarrow{BA}=(2,0,0)$，二面角的余弦值 $\cos\theta = \frac{|\vec{m}\cdot\overrightarrow{BA}|}{|\vec{m}||\overrightarrow{BA}|} = \frac{6}{2\sqrt{2a^2-2a+14}} = \frac{3}{\sqrt{2a^2-2a+14}}$，当 $a=\frac{1}{2}$ 时，分母最小，$\cos\theta$ 最大，$\sin\theta$ 最小，此时 $B_1D = a = \frac{1}{2}$．
\end{solution}
\end{question}

\begin{question}[points=12]
抛物线 $C$ 的顶点为坐标原点 $O$，焦点在 $x$ 轴上，直线 $l: x=1$ 交 $C$ 于 $P, Q$ 两点，且 $OP \perp OQ$．已知点 $M(2,0)$，且 $\odot M$ 与 $l$ 相切．（1）求 $C$，$\odot M$ 的方程；（2）设 $A_1, A_2, A_3$ 是 $C$ 上的三个点，直线 $A_1A_2$，$A_1A_3$ 均与 $\odot M$ 相切．判断直线 $A_2A_3$ 与 $\odot M$ 的位置关系，并说明理由．
\begin{solution}
\textbf{答案：}（1）$C: y^2 = x$，$\odot M: (x-2)^2 + y^2 = 1$；（2）相切，理由见解析．

\textbf{分析：}（1）根据已知抛物线与 $x=1$ 相交，可得出抛物线开口向右，设出标准方程，再利用对称性设出 $P, Q$ 坐标，由 $OP \perp OQ$，即可求出 $p$；由圆 $M$ 与直线 $x=1$ 相切，求出半径；（2）先考虑 $A_1A_2$ 斜率不存在，根据对称性，即可得出结论；若斜率存在，由 $A_1, A_2, A_3$ 三点在抛物线上，将直线 $A_1A_2, A_1A_3, A_2A_3$ 斜率分别用纵坐标表示，再由 $A_1A_2, A_1A_3$ 与圆 $M$ 相切，得出 $y_2+y_3, y_2y_3$ 与 $y_1$ 的关系，最后求出 $M$ 到直线 $A_2A_3$ 的距离，即可得出结论．

\textbf{详解：}（1）设抛物线 $C: y^2 = 2px$（$p>0$），由对称性设 $P(1, \sqrt{2p}), Q(1, -\sqrt{2p})$，$OP \perp OQ$ 得 $1-2p=0$，所以 $p=\frac{1}{2}$，故 $C: y^2 = x$．$\odot M$ 与 $x=1$ 相切，半径 $r=1$，故 $\odot M: (x-2)^2 + y^2 = 1$．（2）设 $A_i(x_i, y_i)$（$i=1,2,3$），则 $x_i = y_i^2$．直线 $A_iA_j$ 方程 $x - (y_i+y_j)y + y_i y_j = 0$．$A_1A_2$ 与 $\odot M$ 相切得 $\frac{|2 + y_1y_2|}{\sqrt{1+(y_1+y_2)^2}} = 1$，整理得 $(y_1^2-1)y_2^2 + 2y_1 y_2 + 3 - y_1^2 = 0$．同理 $A_1A_3$ 得 $(y_1^2-1)y_3^2 + 2y_1 y_3 + 3 - y_1^2 = 0$．故 $y_2, y_3$ 是方程 $(y_1^2-1)y^2 + 2y_1 y + 3 - y_1^2 = 0$ 的两根，$y_2+y_3 = -\frac{2y_1}{y_1^2-1}$，$y_2y_3 = \frac{3-y_1^2}{y_1^2-1}$．则 $M$ 到直线 $A_2A_3$ 的距离 $d = \frac{|2 + y_2y_3|}{\sqrt{1+(y_2+y_3)^2}} = \frac{\left|2+\frac{3-y_1^2}{y_1^2-1}\right|}{\sqrt{1+\left(\frac{2y_1}{y_1^2-1}\right)^2}} = 1$，所以 $A_2A_3$ 与 $\odot M$ 相切．
\end{solution}
\end{question}

\begin{question}[points=12]
已知 $a>0$ 且 $a\neq 1$，函数 $f(x) = \frac{x^a}{a^x}$（$x>0$）．（1）当 $a=2$ 时，求 $f(x)$ 的单调区间；（2）若曲线 $y=f(x)$ 与直线 $y=1$ 有且仅有两个交点，求 $a$ 的取值范围．
\begin{solution}
\textbf{答案：}（1）单调递增区间 $\left(0, \frac{2}{\ln 2}\right]$，单调递减区间 $\left[\frac{2}{\ln 2}, +\infty\right)$；（2）$(1,e) \cup (e, +\infty)$

\textbf{分析：}（1）求得函数的导函数，利用导函数的正负与函数的单调性的关系即可得到函数的单调性；（2）利用指数对数的运算法则，可以将曲线 $y=f(x)$ 与直线 $y=1$ 有且仅有两个交点等价转化为方程 $\frac{\ln x}{x} = \frac{\ln a}{a}$ 有两个不同的实数根，即曲线 $y=g(x)$ 与直线 $y=\frac{\ln a}{a}$ 有两个交点，利用导函数研究 $g(x)$ 的单调性，并结合 $g(x)$ 的正负，零点和极限值分析图象，进而得到 $0<\frac{\ln a}{a}<\frac{1}{e}$，即 $a>1$ 且 $a\neq e$．

\textbf{详解：}（1）$a=2$ 时，$f(x)=\frac{x^2}{2^x}$，$f'(x)=\frac{x\cdot2^x(2-x\ln2)}{4^x}=\frac{x(2-x\ln2)}{2^x}$，由 $f'(x)=0$ 得 $x=\frac{2}{\ln2}$，所以 $f(x)$ 在 $\left(0,\frac{2}{\ln2}\right]$ 单调递增，在 $\left[\frac{2}{\ln2}, +\infty\right)$ 单调递减．（2）$f(x)=1$ 等价于 $a^x = x^a$，即 $x\ln a = a\ln x$，即 $\frac{\ln x}{x} = \frac{\ln a}{a}$．设 $g(x)=\frac{\ln x}{x}$，则 $g'(x)=\frac{1-\ln x}{x^2}$，$g(x)$ 在 $(0,e]$ 单调递增，在 $[e,+\infty)$ 单调递减，最大值为 $g(e)=\frac{1}{e}$．曲线 $y=f(x)$ 与 $y=1$ 有两个交点等价于 $y=g(x)$ 与 $y=\frac{\ln a}{a}$ 有两个交点，故 $0<\frac{\ln a}{a}<\frac{1}{e}$．函数 $g(a)$ 在 $(0,e)$ 单调递增，$(e,+\infty)$ 单调递减，且 $g(1)=0$，$g(e)=1/e$，结合图象得 $a\in(1,e)\cup(e,+\infty)$．
\end{solution}
\end{question}

\section*{选修4-4：坐标系与参数方程}
请考生在第22、23题中任选一题作答．如果多做，则按所做的第一题计分．本小题10分．

\begin{question}[points=10]
在直角坐标系 $xOy$ 中，以坐标原点为极点，$x$ 轴正半轴为极轴建立极坐标系，曲线 $C$ 的极坐标方程为 $\rho = 2\sqrt{2}\cos\theta$．（1）将 $C$ 的极坐标方程化为直角坐标方程；（2）设点 $A$ 的直角坐标为 $(1,0)$，$M$ 为 $C$ 上的动点，点 $P$ 满足 $\overrightarrow{AP} = \sqrt{2}\overrightarrow{AM}$，写出 $P$ 的轨迹 $C_1$ 的参数方程，并判断 $C$ 与 $C_1$ 是否有公共点．
\begin{solution}
\textbf{答案：}（1）$(x-\sqrt{2})^2+y^2=2$；（2）$P$ 的轨迹 $C_1$ 的参数方程为 $\begin{cases} x = 3-\sqrt{2}+\sqrt{2}\cos\theta \\ y = \sqrt{2}\sin\theta \end{cases}$（$\theta$ 为参数），$C$ 与 $C_1$ 没有公共点．

\textbf{分析：}（1）将曲线 $C$ 的极坐标方程化为 $\rho^2 = 2\sqrt{2}\rho\cos\theta$，将 $x=\rho\cos\theta, y=\rho\sin\theta$ 代入可得；（2）设 $P(x,y)$，设 $M(\sqrt{2}+\sqrt{2}\cos\theta, \sqrt{2}\sin\theta)$，根据向量关系即可求得 $P$ 的轨迹 $C_1$ 的参数方程，求出两圆圆心距，和半径之差比较可得．

\textbf{详解：}（1）由 $\rho=2\sqrt{2}\cos\theta$ 得 $\rho^2=2\sqrt{2}\rho\cos\theta$，即 $x^2+y^2=2\sqrt{2}x$，故 $(x-\sqrt{2})^2+y^2=2$．（2）设 $M(\sqrt{2}+\sqrt{2}\cos\theta, \sqrt{2}\sin\theta)$，$P(x,y)$，由 $\overrightarrow{AP}=\sqrt{2}\overrightarrow{AM}$ 得 $(x-1,y)=\sqrt{2}(\sqrt{2}+\sqrt{2}\cos\theta-1, \sqrt{2}\sin\theta)$，即 $\begin{cases} x = 3-\sqrt{2}+\sqrt{2}\cos\theta \\ y = \sqrt{2}\sin\theta \end{cases}$，为 $C_1$ 的参数方程．$C$ 的圆心 $C(\sqrt{2},0)$，半径 $r=\sqrt{2}$；$C_1$ 的圆心 $C_1(3-\sqrt{2},0)$，半径 $R=\sqrt{2}$，圆心距 $|CC_1| = 3-2\sqrt{2} \approx 0.172$，$R-r = 0$，故两圆内含，无公共点．
\end{solution}
\end{question}

\section*{选修4-5：不等式选讲}
请考生在第22、23题中任选一题作答．如果多做，则按所做的第一题计分．本小题10分．

\begin{question}[points=10]
已知函数 $f(x) = |x-2|$，$g(x) = |2x+3| - |2x-1|$．（1）画出 $y=f(x)$ 和 $y=g(x)$ 的图像；（2）若 $f(x+a) \ge g(x)$，求 $a$ 的取值范围．

\begin{center}
  \includegraphics[width=.38\linewidth]{question-23-1.jpg}
\end{center}
\begin{solution}
\textbf{答案：}（1）图像见解析；（2）$a \ge \frac{11}{2}$

\textbf{分析：}（1）分段去绝对值即可画出图像；（2）根据函数图像数形结合可得需将 $y=f(x)$ 向左平移可满足条件，求得 $y=f(x+a)$ 过 $A(\frac{1}{2},4)$ 时 $a$ 的值可求．

\textbf{详解：}（1）$f(x)=|x-2| = \begin{cases} 2-x, & x<2 \\ x-2, & x\ge 2 \end{cases}$，$g(x)=|2x+3|-|2x-1| = \begin{cases} -4, & x<-\frac{3}{2} \\ 4x+2, & -\frac{3}{2}\le x < \frac{1}{2} \\ 4, & x\ge \frac{1}{2} \end{cases}$，图像略．（2）$f(x+a)=|x+a-2|$，由图像知需将 $f(x)$ 向左平移，即 $a>0$．当 $f(x+a)$ 过点 $A(\frac{1}{2},4)$ 时，$|\frac{1}{2}+a-2|=4$，解得 $a=\frac{11}{2}$ 或 $a=-\frac{5}{2}$（舍去），所以 $a\ge\frac{11}{2}$．
\end{solution}
\end{question}

